\documentclass[11pt]{article}
\usepackage[utf8]{inputenc}
\usepackage[T1]{fontenc}
\usepackage{lmodern}
\usepackage[margin=1in]{geometry}
\usepackage{amsmath,amssymb}
\usepackage{textcomp}
\usepackage[dvipsnames]{xcolor}
\usepackage{enumitem}
\usepackage{parskip}
\usepackage{titlesec}
\usepackage[colorlinks=true,linkcolor=blue,citecolor=blue,urlcolor=blue]{hyperref}
\usepackage{natbib}

\titleformat{\section}{\large\bfseries}{\thesection}{1em}{}
\titleformat{\subsection}{\normalsize\bfseries}{\thesubsection}{1em}{}
\titleformat{\subsubsection}{\normalsize\bfseries\itshape}{}{0em}{}

\title{\textbf{The Trace Is Not the Truth:\\ Measuring Attribution Blind Spots under Schema-Valid Tool Corruption}}
\author{}
\date{}

\newcommand{\ph}{\textit{[placeholder --- fill in after experiments are run]}}

\begin{document}
\maketitle
\vspace{-2em}

\begin{quote}
\small\textit{Assembled from the six confirmed stage outputs of an \texttt{auto-research-discussion} run (seed topic B2.5), post research-taste audit (\textcolor{ForestGreen}{\textbf{PASS}}). Revised 2026-08-21 per \texttt{review\_2026-08.md}: conceded-and-positioned against The Misattribution Gap, Who\&When Pro, Don't Blindly Trust It, and the 2026 attribution wave; pre-registered run-budget/descope ladder added. Status: \textbf{specification / paper draft --- no experiments have been run.} All numbers below are pre-registered design targets, not results.}
\end{quote}

\begin{abstract}
Automated failure attribution for LLM agents --- deciding whether a failed trajectory should be blamed on the plan, a tool, or the execution layer --- is being industrialized to drive real interventions: retraining planners, patching tools, curating data. Every attribution method proposed so far, however, treats the trajectory log it reads as trustworthy evidence, so a tool call that returns a semantically wrong but syntactically valid observation is invisible to the very validity checks these pipelines run, and its corrupted content is then reused as evidence against the plan and execution layers that consumed it. We build a schema-preserving corruption injector gated by three independent validity checks, verify each corruption's causal decisiveness by counterfactual twin-replay, and run existing LLM-judge, trained-tracer, and runtime-monitor attribution methods --- blinded to injection metadata --- against matched overt-error and agent-side-fault controls. A one-sentence trust-license prompt intervention probes whether the resulting bias is a removable trust assumption or a deeper reasoning limitation, and a pre-registered, dual-annotated mapping rubric scores every method's native output against a shared plan/tool/execution ground truth. We report the misattribution rate, its direction, and the mechanism-probe outcome for current attribution methods, and what this implies for automated debugging pipelines that treat the trace as ground truth.
\end{abstract}

\section{Introduction}

Multi-agent and tool-using LLM systems fail often enough, and in different enough places, that manually reading the trajectory log to decide whose fault it was has become a bottleneck. Automated failure attribution --- pointing an LLM judge or a trained tracer at the log and asking it to name the responsible agent and step --- is being industrialized precisely to remove that bottleneck: the Who\&When benchmark \citep{WhoWhenGazlTYxZss} motivates its own release by noting that manual attribution is ``prohibitively expensive and inconsistent,'' and AgenTracer \citep{AgenTracerl05DseqvuD} frames its tracer's output as ``the missing interface between failure observation and system improvement.'' In both framings, the attribution verdict is not academic: it is meant to tell a team which component to retrain, patch, or retry.

Yet a central assumption behind this automation remains untested on the one channel every attribution method reads at every step: each of these methods consumes the trajectory log's tool observations as trustworthy evidence of what happened, with no step that asks whether those observations might themselves be wrong. The memory-poisoning literature has recently documented the analogous failure for a \emph{different} evidence channel --- adversarially poisoned persistent memory that defenders consistently misread as model failure \citep{MisattributionGap260522842} --- but the transient tool-observation channel, corrupted not by an attack but by a plainly wrong value, has no corresponding measurement. A tool call that returns a value which is syntactically valid, parses cleanly, and passes every schema check the pipeline runs --- but is simply the wrong number, the wrong entity, or a stale reading --- leaves no surface trace for a validity check to catch. Worse, the same trajectory log that constitutes the failure's evidence is produced by the very tool call under suspicion, so a diagnostic procedure that reads the log at face value inherits the corruption into its own premises before it starts reasoning about whose fault the failure is. If attribution methods systematically relabel these tool-caused failures as plan or execution failures, every downstream fix they drive is aimed at a component that was never broken, while the actual defect --- trusting well-formed tool output as ground truth --- goes unaddressed and unmonitored.

The nearest existing work approaches this gap from adjacent but distinct angles. Who\&When \citep{WhoWhenGazlTYxZss} formalizes attribution as identifying the agent-step pair whose correction flips a failed trajectory to success, and evaluates LLM-judge attribution under several context-management strategies --- but its decisive-error definition is exercised on naturally occurring agent mistakes, never on a condition where the log's own tool observations are corrupted yet well-formed. AgenTracer \citep{AgenTracerl05DseqvuD} trains a dedicated tracer using counterfactual replay and programmatic fault injection, but its fault universe is agent actions, not tool returns --- a tracer trained to recognize a bad action has, by construction, never seen a trajectory whose decisive error is a lying observation. Who\&When Pro \citep{WhoWhenPro260709996} scales injection-built attribution evaluation to 12,326 trajectories with golden labels assigned by convention at injection time, yet nothing in its construction certifies that an injected fault is surface-undetectable, and it measures who attributes well over logs presumed faithful, not what happens to attribution when the log itself lies. A concurrent causal-inference approach, Causal Agent Replay \citep{CausalAgentReplay260608275}, shares our reliance on a \texttt{do\_observation}-style counterfactual substitution, but proposes and validates a replacement attributor against synthetic causal models with mocked tools, rather than measuring whether existing, deployed methods misattribute blame under a verified, schema-valid corruption of a real tool observation --- the two efforts share machinery, not a research question. Closest in spirit, the Misattribution Gap \citep{MisattributionGap260522842} documents misattribution under a corrupted evidence channel --- memory-poisoning attacks that attribution and safety tooling consistently misread as model failure --- but its channel is a persistent shared memory store seeded with an adversarial policy document, and its diagnosis targets are safety classifiers and defender remediation, not the plan/tool/execution attribution families measured here. Don't Blindly Trust It \citep{DontBlindlyTrust260621409} manipulates the tool-feedback channel directly, contrasting faithful, misleading, and absent observations in a matched agent loop, but measures the agent's task performance under bad evidence, not the diagnosis of the agent --- misleading observations are known to break agents; whether they also break the \emph{attribution of blame for} agents is this paper's question. Security-oriented work on the tool-return channel, such as ChatInject and MCP-SafetyBench, shows that tool output can be actively corrupted, but measures whether the resulting attack succeeds, not whether any downstream diagnosis correctly localizes it --- and their corruptions are instruction-bearing payloads designed to redirect behavior, a different threat model from a plainly wrong value in an otherwise ordinary field. ResiliBench models tool unreliability as an overt signal --- timeouts, invalid input, visible operational failure --- that an agent must adapt to, leaving the silent case of a well-formed but wrong return outside its scope. And a lightweight runtime supervisor such as SupervisorAgent triggers on explicit error occurrence, repetitive behavior, or excessive observation length --- precisely the signals a schema-valid corruption does not emit.

Turning this gap into a measurable claim surfaces its own difficulties. Manufacturing a corrupted tool observation whose ground truth is verifiably the decisive cause of failure --- not merely correlated with it --- requires ruling out that the agent simply recovered, or that the corrupted field was never load-bearing in the first place. That same corruption must also leave no surface signal for a validity check to catch, which is a genuinely separate property from being decisive: a corruption can be undetectable yet incidental to the failure, or detectable in principle yet still the actual cause, and conflating the two would make any measured bias an artifact of the construction rather than a property of the methods being tested. Once a corrupted trajectory exists, separating a genuinely broken plan from a plan that was a locally reasonable response to bad information requires a counterfactual --- what would this step have looked like given the correct observation --- that no attribution method under test has access to, only the experimenter's own injected ground truth does. And because different attribution methods report their verdicts in different native vocabularies --- an agent name, a free-text rationale, a discrete monitor flag --- scoring them against one shared notion of ``tool,'' ``plan,'' or ``execution'' requires a mapping choice that is itself capable of manufacturing or hiding a bias if left unprincipled.

Our approach turns each of these obstacles into a specific design choice rather than a hedge. A schema-preserving injector mutates one tool observation's content while three independent checks --- schema re-validation, parser acceptance, and a contextual-plausibility screen --- jointly certify that the corruption is legal and unremarkable, so undetectability is verified rather than assumed. A counterfactual twin-replay procedure keeps only trajectories where undoing just the corrupted observation flips the outcome, supplying the decisive-cause guarantee that undetectability alone cannot. Existing attribution methods are then run blinded to injection metadata against two matched controls --- overt tool errors sharing the same root cause but a visible signal, and the identical wrong value injected into the agent's own action instead --- so that any measured bias can be attributed specifically to the observation channel rather than to the generic difficulty of spotting a subtly wrong value. A one-sentence trust-license addition to the judge prompt then tests whether a bias, once found, is a removable assumption or a deeper limitation, and a mapping rubric fixed before any result is seen --- with its own reliability reported --- keeps the plan/tool/execution scoring itself from being a hidden source of the finding.

We evaluate this design over tool-calling tasks from AgentDojo and long-horizon assistant tasks from GAIA, using Who\&When's human-annotated logs as an independently sourced control of genuine plan failures, against a naive majority-class attributor, matched overt-error trajectories, and a matched agent-side-fault arm.

In summary, we contribute:
\begin{itemize}[leftmargin=1.4em]
\item An identification and formalization of a structural asymmetry in automated failure attribution: validity checks certify a tool observation's form, while the diagnosis itself consumes the observation's content as an unquestioned premise, making a schema-valid lie both invisible to detection and actively reused as evidence against the layers that consumed it.
\item A counterfactual-replay-verified corruption testbed that separates causal decisiveness from surface undetectability as two independently checkable properties of an injected fault, rather than conflating them into a single validity check.
\item A matched-control protocol --- overt tool errors and an agent-side-fault arm sharing the identical wrong value at a different causal locus --- that isolates whether any measured attribution bias is specific to the tool-observation channel.
\item A minimal trust-license prompt intervention that tests whether a measured bias is a removable structural trust assumption or a deeper counterfactual-reasoning limitation, yielding a concrete, falsifiable design recommendation either way.
\item An empirical characterization of the direction of misattribution and its dependence on corruption type, locus, and subtlety --- not just whether current attribution methods fail under lying evidence, but how the failure is distributed.
\end{itemize}

\section{Problem Definition}

\noindent\textbf{Formal setup.} Consider a trajectory $\tau = (q, (a_1,o_1), \dots, (a_T,o_T))$ pairing a task $q$ with a sequence of agent actions $a_t$ and tool observations $o_t$, scored by a programmatic checker $V(\tau) \in \{0,1\}$. A corruption operator $\mu$ maps a chosen observation $o_t$ to a mutated $o'_t$ while preserving schema and parse validity --- the mutated value type-checks and is accepted by the framework's own parser exactly as the original would be. A twin-replay operator $\mathrm{replay}(\tau, t, o)$ re-executes $\tau$ from step $t$ with observation $o$ substituted, holding seed and decoding temperature fixed where the framework allows. An attribution method $A$ maps a (possibly corrupted) trajectory to a responsible layer and step, $A(\tau) = (\hat{\ell}, \hat{t})$ with $\hat{\ell} \in \{\text{plan}, \text{tool}, \text{execution}\}$, via a method-specific mapping rubric $R_A$ that converts the method's native output into this shared label space. The object of study is whether, for trajectories whose ground truth is $\ell = \text{tool}$ at the injection step $t$, $A$ instead reports $\hat{\ell} \neq \text{tool}$ at a rate exceeding the same method's error rate on matched trajectories where the tool fault is overt rather than schema-valid.

\noindent\textbf{Why this is hard.} Fixing $\mu(o_t) = o'_t$ does not by itself establish that $o'_t$ is the reason $V(\tau) = 0$: the corrupted run could fail (or succeed) for reasons unrelated to the mutated field, so a ground-truth label of $\ell = \text{tool}$ at step $t$ is only defensible if $V(\tau_{\text{corrupted}}) = 0$ and $V(\mathrm{replay}(\tau_{\text{corrupted}}, t, o_t)) = 1$ --- a twin-replay condition, not a property of $\mu$ alone. Independently, $\mu$ must produce an $o'_t$ indistinguishable from a valid observation by any surface check available to $A$, a property orthogonal to the twin-replay condition: a corruption can satisfy one without the other. Even once both hold, $R_A$ is itself a design choice --- different attribution methods report $\hat{\ell}$ in incompatible native vocabularies, and a mapping generous or strict enough to change $\hat{\ell}$'s value after the fact would make a measured misattribution rate an artifact of $R_A$ rather than of $A$. And distinguishing ``$A$ blamed the plan because the plan was genuinely defective'' from ``$A$ blamed the plan because that was a reasonable read of a corrupted $o_t$'' requires the counterfactual outcome under the true $o_t$ --- exactly what $\mathrm{replay}(\tau, t, o_t)$ supplies to the experimenter but withholds from $A$ itself, since $A$ must diagnose blind, as any deployed method would.

\section{Method Overview}

At a conceptual level, the study is a single measurement pipeline with five cooperating parts, each converting one of the difficulties above into a checkable property rather than an assumption. A trajectory and a target tool observation enter the pipeline; a corruption-and-verification stage turns a candidate mutation into a ground-truth-labeled treatment case only if it clears both a legality bar and a causal-decisiveness bar; existing attribution methods are then run, blinded, against that treatment case and two matched controls; a trust-license variant of each LLM-judge method probes the mechanism behind any measured bias; and a pre-registered scoring rubric turns every method's heterogeneous native output into one comparable verdict.

The \textbf{schema-preserving corruption injector} takes a candidate tool observation and a mutation operator and produces a corrupted variant, passing it forward only if it clears three independent legality gates --- the surface-undetectability half of ground-truth construction. The \textbf{counterfactual twin-replay oracle} takes a legal candidate and confirms, by re-executing the trajectory with the original observation restored, that the corruption is the decisive cause of failure --- the causal-decisiveness half. The \textbf{blinded attribution harness} takes the resulting verified treatment set, together with matched overt-error and agent-side-fault control sets, and runs every attribution method under test over all three, blind to which set a given trajectory came from. The \textbf{trust-license probe} re-runs each LLM-judge method with one added prompt sentence licensing distrust of tool observations, isolating whether any measured bias is a removable trust assumption. The \textbf{pre-registered mapping rubric} converts each method's native output into the shared plan/tool/execution label before any result is seen, with its own inter-annotator reliability reported alongside every downstream number.

\section{Methodology}

\subsection{Schema-Preserving Corruption Injector}
\textbf{Challenge.} A surface-detectable corruption would only measure whether attribution methods bother to validate their input, not whether they can diagnose a well-formed lie.

\textbf{Insight.} No single validity check can certify invisibility, since each check tests a different thing --- type structure, runtime parsing, contextual plausibility --- so passing all three simultaneously is what makes a corruption both legal and unremarkable.

\textbf{Method.} A library of typed mutation operators --- numeric perturbation, same-type entity substitution, boolean flip, list truncation/reorder, and stale-value substitution --- is applied to one tool return per trajectory at a controlled step index. Each candidate is kept only if it (a) re-validates against the tool's declared schema/return type, (b) is accepted by the framework's native parser without raising an error, and (c) is judged ``not self-evidently absurd in context'' by a plausibility-screen model disjoint from any judge or agent backbone family.

\subsection{Counterfactual Twin-Replay Oracle}
\textbf{Challenge.} Undetectability alone does not establish that a corruption caused the failure --- the corrupted run could fail for an unrelated reason, or the agent could simply recover.

\textbf{Insight.} Causal decisiveness is a separate property from surface stealth; only re-executing the trajectory with the original observation restored rules out an undetectable-but-incidental corruption.

\textbf{Method.} Each legality-gate-passing candidate is paired with a twin replay that substitutes the original observation back in at the same step, holding seed and temperature fixed where the framework allows. A candidate is promoted to the verified treatment set only if the corrupted run fails the task's programmatic checker and the twin succeeds --- over-generating 3--5x to absorb the yield loss from candidates that fail this test.

\subsection{Blinded Attribution Harness and Trust-License Probe}
\textbf{Challenge.} A measured misattribution rate could reflect a hard capability ceiling rather than a fixable default of treating the trace as trustworthy evidence, and the two explanations call for different remedies.

\textbf{Insight.} If a one-sentence license to distrust evidence a method already has access to measurably changes its verdict, the bias is better explained by a structural trust default than by a capability limit, since capability did not change --- only the license did.

\textbf{Method.} Every attribution method under test --- Who\&When's \citep{WhoWhenGazlTYxZss} LLM-judge variants, an AgenTracer-style \citep{AgenTracerl05DseqvuD} trained tracer, and a SupervisorAgent-style \citep{SupervisorAgentpzFhtpkabh} runtime monitor replayed post hoc --- is run blinded to injection metadata over the verified treatment set and both control sets. For each LLM-judge variant, a second run adds exactly one prompt sentence stating that tool observations may be silently wrong and should be independently sanity-checked; the paired difference in misattribution rate, with and without this sentence, is the mechanism readout.

\subsection{Matched Control Design}
\textbf{Challenge.} A raw misattribution rate on corrupted trajectories cannot say whether the bias is specific to tool observations or just reflects the general difficulty of noticing any subtly wrong value.

\textbf{Insight.} The experimenter's own replay-verified ground truth supplies exactly the counterfactual comparison no attribution method under test has access to, turning ``the plan step looks defensible given what the agent saw'' from an unfalsifiable excuse into a measured quantity.

\textbf{Method.} Two control sets are constructed on the same tasks and step loci as the treatment set: an overt-tool-error arm (exceptions, timeouts, empty returns) sharing the tool-layer root cause but emitting a visible signal, and an agent-side-fault arm that injects the identical wrong value into the agent's own action rather than the tool's observation, matched for semantic delta and subtlety. Misattribution is compared across all three arms, and separately cross-tabulated by which layer was blamed (plan vs.\ execution) as a function of corruption operator family, injection locus (early/mid/late tercile), and subtlety.

\subsection{Pre-Registered Mapping Rubric}
\textbf{Challenge.} Attribution methods report verdicts in incompatible native vocabularies, and an unprincipled mapping onto a shared plan/tool/execution ontology could manufacture or hide a bias after the fact.

\textbf{Insight.} Fixing the mapping before any result is seen removes the experimenter's own opportunity to choose a generous or strict mapping once the answer it would produce is already known.

\textbf{Method.} One deterministic mapping rubric per attribution method is written before any experimental result is examined, converting that method's native output (an agent name, a free-text rationale, a discrete flag) into $\{\text{plan}, \text{tool}, \text{execution}\}$. The rubric is validated by dual annotation on a held-out sample, with Cohen's $\kappa$ reported per method as a precondition for trusting any downstream misattribution number; a sensitivity analysis excludes rubric-disputed trajectories from the headline comparison.

\section{Experiments}

\subsection{Experimental Setup}

\textbf{Datasets.} AgentDojo \citep{AgentDojo240613352} --- real tool-calling tasks (email, banking, travel, Slack) with programmatic checkers and structured, injectable tool returns; the pilot-stage domain, cheapest to instrument the injector and replay oracle against. GAIA \citep{GAIA231112983} --- long-horizon, multi-step assistant tasks where a single early observation is load-bearing many steps later; used specifically to stress-test the locus factor (early/mid/late injection tercile) and the type-dependence of misattribution direction under longer causal chains than AgentDojo's shorter tasks provide. Who\&When's \citep{WhoWhenGazlTYxZss} human-annotated failure logs --- a source of genuine plan-failure controls with independently verified clean observations, avoiding the circularity of self-labeling ``genuine'' plan failures from synthetic data alone.

\textbf{Comparators} (sanity-check floors, not competitors to beat --- this is a measurement study, and no research question is framed as ``our procedure beats baseline X''). A constant ``always-predict-plan'' attributor floors the comparison: any method whose accuracy on the corrupted set does not beat naive majority-guessing is not performing attribution on corrupted evidence at all. The overt-tool-error arm establishes that methods can localize tool faults when the fault emits a visible signal --- the necessary comparator for every corrupted-observation claim. The matched agent-side-fault arm --- the identical wrong value, injected into the agent's action instead of the tool's observation --- isolates whether any bias is specific to the observation channel.

\textbf{Placebo validity gate (checked before any research question is scored).} Neither the overt-error arm nor the constant attributor carries a null prediction, so neither can certify that the pipeline itself is artifact-free. Two null-by-construction components are evaluated first: a \emph{sham-injection arm} that runs the identity mutation ($o'_t = o_t$) through the identical injection-replay-serialization machinery, requiring each method's attribution distribution on sham-processed versus untouched twins to be statistically indistinguishable (equivalence margin $\pm 3$ percentage points); and a \emph{label-shuffle placebo} that permutes the semantic-corruption/overt-error arm labels within matched task-step pairs 1{,}000 times, requiring the permuted $\Delta$-trust gap to be centered on zero. A non-null result on either component invalidates the run in full --- no research question below may be scored from an invalid run, per the general rule that a pipeline capable of manufacturing its own headline effect must be ruled out before that effect is trusted.

\textbf{Premise-collapse and type-independence branches.} If methods cannot reliably localize even overt, signal-emitting tool faults ($m_A(D_{\text{overt}}) > 40\%$ for a majority of methods), the schema-validity distinction this study is built to measure is not yet interpretable; in that cell the finding becomes a competence audit anchored on the constant-attributor floor, and RQ5--RQ6 below are reported descriptively only, not scored as pass/fail. Independently, if no corruption-type, locus, or subtlety factor predicts blame direction (RQ6), the honest finding is that the blind spot is indiscriminate across corruption families rather than type-targeted --- a simpler, still-reportable failure geometry, distinct from RQ6's direction claim itself surviving or failing.

\textbf{Human-annotator reference arm (descriptive, pre-registered).} A stratified 10\% subsample of the verified treatment set is additionally annotated by human raters blind to injection metadata, using the same plan/tool/execution ontology. This arm carries no pass/fail bar; it is registered up front so that neither direction of the human result can be cherry-picked after the fact. If humans are also fooled, human-annotated attribution ground truths --- the labels and taxonomies existing benchmarks are built on --- are themselves exposed to the corruption channel this study measures, reported as a scope statement on the sampled strata; if humans are not fooled, the human rate becomes the repair upper bound against which the trust-license reduction (RQ3--RQ4) is compared.

\textbf{Run budget and descope ladder (pre-registered).} At scale-up targets (three arms $\times$ $\geq 300$ verified trajectories per arm per framework $\times$ 3 frameworks), the verified corpus alone comprises ${\approx}2{,}700$ trajectories; at the RQ2 verified-yield floor of 20\% on load-bearing candidates, constructing the semantic-corruption arm implies ${\approx}4{,}500$ gate-passing candidates per framework --- ${\approx}13{,}500$ total, each requiring a corrupted run plus a twin replay, i.e.\ on the order of 27{,}000 constructions before attribution scoring (5 methods over all arms, plus paired trust-license re-runs of the three judge variants). If pilot-stage cost accounting shows this exceeding budget, the pre-registered descope order is: (1) drop the third framework, retaining AgentDojo and GAIA with the full arm structure; (2) demote the GAIA locus-stratification factor to a descriptive readout; (3) reduce per-arm $n$ toward the RQ3/RQ6 power-band minima, reporting any cell that falls below its band as INCONCLUSIVE rather than silently under-powered. No descope step may drop an arm: the matched-control contrast is the paper's identification strategy, and losing it converts the study into an uninterpretable single-arm rate.

\textbf{Metrics.} Layer-attribution accuracy and its complement, the misattribution rate; the $\Delta$-trust gap (misattribution rate on schema-valid corruption minus on matched overt-error controls); the direction share among misattributed cases (proportion blamed to plan vs.\ execution); step-localization accuracy (exact and $\pm 1$); the control false-positive rate (genuine plan failures wrongly blamed on the tool layer); and Cohen's $\kappa$ on the mapping rubric.

\textbf{Research questions:}
\begin{enumerate}[leftmargin=1.6em]
\item RQ1 --- Does the automated legality screen approximate human undetectability?
\item RQ2 --- Does load-bearing-field targeting drive verified corruption yield?
\item RQ3 --- Does a trust-license prompt reduce misattribution?
\item RQ4 --- Is any trust-license reduction partial rather than complete?
\item RQ5 --- Is the misattribution bias specific to the tool-observation channel?
\item RQ6 --- Does misattribution skew toward the plan layer, and does the skew depend on corruption type?
\item RQ7 --- Is the plan/tool/execution mapping rubric reliable?
\end{enumerate}

\subsubsection{RQ1: Does the automated legality screen approximate human undetectability?}
\textbf{Setup.} Draw a stratified sample of 150 corrupted observations that pass all three automated legality gates (30 per operator family), embedded in their surrounding trajectory context, alongside a calibration set of 50 candidates the injector rejected at the plausibility gate. Present both sets to an independent human rater blind to which observations were corrupted or rejected.

\textbf{Expected outcome.} At least 90\% of gate-passing corruptions are independently rated plausible-in-context, versus a substantially lower rate (target: below 60\%) for the rejected-candidate calibration set.

\textbf{Falsification condition.} Refuted if the gate-passing plausibility rate falls below 90\%, or if it is not substantially higher than the rejected-candidate calibration set's rate --- either would mean the automated screen is a weak proxy for human undetectability, and any downstream misattribution claim would need to be re-examined for whether it is measuring diagnosis under a genuinely invisible corruption or under one a careful reader would already flag.

\textbf{Hypothesis (plain language).} The automated legality gates are a reasonably faithful proxy for whether a human reading the trajectory would notice anything wrong.

\textbf{Metrics.} Human-rated plausibility rate on gate-passing vs.\ rejected-candidate sets.

\textbf{Result.} \ph

\subsubsection{RQ2: Does load-bearing-field targeting drive verified corruption yield?}
\textbf{Setup.} Attempt corruption injection on load-bearing fields (fields referenced by the agent's own subsequent step in the uncorrupted run, identified by string/AST trace matching) versus non-load-bearing fields, matched for operator family and task domain; compute twin-replay-verified yield separately for each group over 100 pilot candidates per group.

\textbf{Expected outcome.} Verified yield on load-bearing-field candidates is at least 20\%, versus a substantially lower yield (target: below 10\%) on non-load-bearing-field candidates.

\textbf{Falsification condition.} Refuted if load-bearing-field yield falls below 20\%, or if the two groups' yields are statistically indistinguishable --- either would mean load-bearing-field targeting is not what drives verified yield, and the construction pipeline would need a different targeting heuristic (or larger over-generation) before the scale-up stage is feasible on budget.

\textbf{Hypothesis (plain language).} Targeting fields the agent's own later steps actually use, not merely corrupting any field that passes the legality gates, is what makes verified ground-truth construction practical.

\textbf{Metrics.} Twin-replay verified yield rate per group.

\textbf{Result.} \ph

\subsubsection{RQ3: Does a trust-license prompt reduce misattribution?}
\textbf{Setup.} Within-method paired design: run each of the three LLM-judge variants twice on the same $\geq 300$ verified corrupted trajectories per framework --- once with, once without, the one-sentence trust-license addition. Because McNemar power depends on the discordant-pair rate rather than on $n$ alone, the reduction clause is scored per judge variant only when at least 72 discordant pairs are observed ($\approx$0.8 power at the pre-registered 5-point reduction); below that, the reduction verdict is pre-declared inconclusive and the paired difference is reported with a confidence interval rather than as a null.

\textbf{Expected outcome.} For at least 2 of 3 judge variants, the license reduces the misattribution rate by at least 5 percentage points absolute (paired McNemar, $\alpha = 0.05$) without a significant false-positive increase on genuine plan-failure controls (non-inferiority margin: 3 percentage points).

\textbf{Falsification condition.} Refuted if the license yields no significant misattribution reduction for a majority of judge variants, or if any reduction is bought with a false-positive increase on plan-failure controls exceeding the 3-point margin --- the former would mean the bias is not simply a removable default, and the latter would mean the license trades one error for another rather than repairing attribution. A significant \emph{increase} in misattribution under the license is pre-registered as a wrong-direction result reported at full prominence --- the license actively harming attribution --- rather than being folded into the no-reduction wording.

\textbf{Hypothesis (plain language).} Judges are not incapable of using distrust information; they default to not applying it unless prompted to.

\textbf{Metrics.} Paired misattribution-rate reduction per judge variant; false-positive rate on plan-failure controls. As a pre-registered descriptive readout with no hypothesis attached, each judge rationale is additionally logged for whether the corrupted observation was flagged as suspicious yet blame was still placed elsewhere --- separating ``never noticed'' from ``noticed but overrode'' at no extra run cost.

\textbf{Result.} \ph

\subsubsection{RQ4: Is any trust-license reduction partial rather than complete?}
\textbf{Setup.} Same paired runs as RQ3; compare post-license misattribution rate against the same method's misattribution rate on matched overt-error controls.

\textbf{Expected outcome.} The license reduces but does not eliminate misattribution --- post-license misattribution rate remains above the method's overt-error-control rate for a majority of judge variants, leaving a majority (target: above 50\%) of the original $\Delta$-trust gap intact.

\textbf{Falsification condition.} Refuted toward ``fully fixed'' if post-license misattribution rate is statistically indistinguishable from the overt-error-control rate for a majority of judge variants; refuted toward ``no effect at all'' if this coincides with RQ3's own falsification (no significant reduction). A result landing between these --- some reduction, but not to the overt-error-control floor --- is the only outcome consistent with the partial-repair hypothesis as stated.

\textbf{Hypothesis (plain language).} Knowing evidence may lie is not the same as knowing which piece of evidence lies; residual misattribution reflects this remaining counterfactual-reasoning gap.

\textbf{Metrics.} Proportion of the original $\Delta$-trust gap remaining after the license, per judge variant.

\textbf{Result.} \ph

\subsubsection{RQ5: Is the misattribution bias specific to the tool-observation channel?}
\textbf{Setup.} Pilot: one single-agent tool-calling framework, 2 backbone judge families, 100 verified trajectories per arm (semantic-corruption, overt-error, agent-side-fault) on AgentDojo. Scale-up: 3 frameworks, $\geq 300$ verified trajectories per arm per framework; the scale-up additionally crosses one judge capability tier --- a weak and a frontier judge within one model family, on one framework --- pre-registering the check that the gap does not vanish at the frontier tier (ruling out the rival reading of a weak-judge artifact that the next model generation fixes). Per-method two-proportion test (McNemar for task-matched pairs) between arms, Holm-corrected across the 5 attribution methods. Because an overt error typically truncates the trajectory at the fault step while a schema-valid corruption lets the agent run on, fault-to-manifestation distance is a named confound of this contrast: the $\Delta$-trust gap is additionally reported stratified by post-fault suffix length, with the pre-committed interpretation that if the gap collapses within suffix-length strata, it is attributed to fault-to-manifestation distance rather than to the observation channel, and channel specificity rests solely on the agent-side-fault contrast.

\textbf{Expected outcome.} The $\Delta$-trust gap (misattribution on semantic corruption minus on the overt-error control) is at least 10 percentage points and significantly positive (Holm-corrected $p < .05$) for a strict majority of the 5 methods --- conditional on those methods clearing the overt-error competence bar ($m_A(D_{\text{overt}}) \leq 40\%$; see the premise-collapse branch above) --- while misattribution on the agent-side-fault arm stays within a few points of the overt-error control. The majority criterion is pre-committed as a fraction --- a strict majority of however many baselines survive the pre-registered smoke test --- so the rule stays defined if a baseline is dropped.

\textbf{Falsification condition.} Refuted if the $\Delta$-trust gap is statistically indistinguishable from zero (or negative) for the majority of methods --- methods would then localize schema-valid tool corruption about as reliably as overt errors. A \emph{significantly negative} gap --- methods better on semantic corruption than on overt errors --- is pre-registered as a wrong-direction finding reported at full prominence in its own right, not folded into the null wording. Refuted in the channel-specificity sub-claim if the agent-side-fault arm's misattribution rate is comparable to or exceeds the tool-observation arm's --- that would mean any bias reflects a general difficulty in spotting subtly wrong values, not something specific to the tool-observation channel, and the ``trust asymmetry'' framing would not survive this test.

\textbf{Hypothesis (plain language).} Attribution methods are not simply bad at spotting subtly wrong values wherever they occur; they are specifically bad at it when the wrong value arrives as a tool observation rather than an agent action.

\textbf{Metrics.} $\Delta$-trust gap per method; agent-side-fault-arm misattribution rate relative to the overt-error control.

\textbf{Result.} \ph

\subsubsection{RQ6: Does misattribution skew toward the plan layer, and does the skew depend on corruption type?}
\textbf{Setup.} Among RQ5's misattributed cases, cross-tabulate blamed layer (plan vs.\ execution) by operator family, injection locus tercile, and subtlety, on AgentDojo and GAIA; exact binomial test of plan-share against 0.5, plus a mixed-effects logistic regression (fixed effects: the three grid factors, with post-fault suffix length included as an additional covariate control; random intercepts: task and framework). \textbf{Power-based scoring bands (corrected from an initial, under-powered $n{=}50$ design):} at the pre-registered target of $\pi_{\text{plan}} = 0.65$, an exact binomial test at $n=50$ has only $\approx$0.51 power, which would misreport a true skew as ``directionally diffuse'' about half the time; the scoring rule is therefore per-cell rather than a single fixed $n$ --- $n \geq 120$ misattributed cases scores the direction claim (power $\approx$0.89 at the critical count 72/120), $50 \leq n < 120$ is pre-declared INCONCLUSIVE rather than scored as a null, and $n < 50$ is reported descriptively only.

\textbf{Expected outcome.} Pooled plan-share among misattributed cases is significantly above 0.5 (target: at least 0.65), and at least one grid factor is a significant predictor of blame direction (likelihood-ratio test, $p < .05$), with the skew toward plan strengthening for subtle, late, numeric/stale-value corruptions.

\textbf{Falsification condition.} Refuted on direction if plan-share does not differ significantly from 0.5 (exact binomial), or skews toward execution instead. Refuted on type-dependence if no grid factor (operator family, locus tercile, subtlety) is a significant predictor in the mixed-effects model --- either result would mean misattribution, when it occurs, is directionally diffuse rather than channeled toward one layer, a distinct and still-reportable failure geometry rather than the predicted skew.

\textbf{Hypothesis (plain language).} When misattribution happens, it is not a coin flip between plan and execution --- plan critique sits immediately downstream of reading the corrupted observation, while execution verdicts key on mechanical signals the corruption does not trip.

\textbf{Metrics.} Plan-share and its confidence interval; mixed-effects model coefficients for the three grid factors.

\textbf{Result.} \ph

\subsubsection{RQ7: Is the plan/tool/execution mapping rubric reliable?}
\textbf{Setup.} Dual annotation of a held-out sample (target: 50 outputs per method across all 5 methods, 250 total) applying the pre-registered rubric before any experimental result is examined; sensitivity re-analysis of the RQ5/RQ6 headline comparisons excluding rubric-disputed trajectories.

\textbf{Expected outcome.} Cohen's $\kappa$ is at least 0.7, pooled and per method, and the sign and significance of the RQ5/RQ6 headline comparisons are unchanged when disputed trajectories are excluded (magnitude may shift).

\textbf{Falsification condition.} Refuted if $\kappa$ falls below 0.7 for one or more methods, or if excluding disputed trajectories flips the sign or significance of the RQ5/RQ6 headline comparisons --- either would mean the mapping rubric itself, not the attribution methods' behavior, is a material driver of the paper's central findings, and those findings would need to be reported as mapping-contingent rather than robust.

\textbf{Hypothesis (plain language).} The plan/tool/execution scoring is itself a reliable instrument, not a hidden source of the measured bias.

\textbf{Metrics.} Cohen's $\kappa$ per method and pooled; headline-result stability under the disputed-trajectory exclusion.

\textbf{Result.} \ph

\vspace{0.5em}
\noindent\textbf{Further evaluation.} With more time, a natural extension is compounding, multi-observation corruption within a single trajectory --- explicitly out of scope here because it changes the identifiability problem from ``is this one injected fault the decisive cause'' to a joint-attribution problem across multiple candidate causes, which would need its own verification design rather than an extension of the twin-replay oracle above.

\section{Related Work}

The closest prior work is Who\&When and AgenTracer: Who\&When already formalizes attribution as identifying a decisive agent-step pair from an LLM judge reading the trajectory log, and AgenTracer already trains a tracer using counterfactual replay and programmatic fault injection --- but neither controls for, nor represents, the case where the log's own tool observations are corrupted yet well-formed, which is precisely the gap this paper is built to measure. Two 2026 works sit closest on either flank and are conceded up front: Who\&When Pro \citep{WhoWhenPro260709996} scales injection-built attribution benchmarking to 12,326 golden-labeled trajectories --- with labels assigned by convention at injection time, no surface-undetectability certification, and no matched channel controls --- and the Misattribution Gap \citep{MisattributionGap260522842} documents misattribution under a corrupted evidence channel for adversarially poisoned persistent memory, a sibling channel and threat model to the non-adversarial, transient tool-observation corruption measured here (positioned in full below). A concurrent causal-inference method, Causal Agent Replay \citep{CausalAgentReplay260608275}, shares our reliance on a same-mechanism observation-substitution intervention but proposes a replacement attributor rather than measuring existing methods' bias --- see below.

\textbf{Automated failure attribution for LLM agents.} Who\&When \citep{WhoWhenGazlTYxZss} formalizes attribution as identifying the decisive agent-step pair whose correction flips a failed trajectory to success, releasing a human-annotated dataset and evaluating LLM-judge attribution under several context-management strategies. AgenTracer \citep{AgenTracerl05DseqvuD} converts decisive-error localization into a trainable prediction problem, generating supervision via counterfactual replay and programmatic fault injection into agent actions, then training a tracer with reinforcement learning. MAST \citep{MAST250313657} derives an empirically grounded taxonomy of recurring multi-agent failure modes from human annotation of naturally occurring traces. Causal Agent Replay \citep{CausalAgentReplay260608275} models an agent run as a structural causal model and attributes blame via \texttt{do(\textperiodcentered)} interventions --- including a \texttt{do\_observation} operation that substitutes a tool result, the same conceptual move as our twin-replay oracle --- but validates only against synthetic causal models with mocked, side-effect-free tools, and does not test whether existing, deployed attribution methods misattribute blame specifically under a verified, schema-valid semantic corruption of a real tool observation; it proposes a new attributor rather than auditing the ones already in use. Who\&When Pro \citep{WhoWhenPro260709996} industrializes this family: a strictly controlled pipeline injects a failure only after exactly replaying a successful prefix, yielding 12,326 failed trajectories with golden labels across 3 modalities and 26 benchmarks --- but the labels are assigned by convention at the moment of injection, with no undetectability certification, no decisiveness round-trip against a matched control, and no condition in which the injected fault is a well-formed lie in the observation channel. A rapidly crowding 2025--26 wave extends the method families further --- CHIEF \citep{CHIEF260223701} converts flat logs into hierarchical causal graphs with counterfactual screening of root causes from propagated symptoms, FALAT \citep{FALAT260600765} frames attribution as dependency-guided search precisely because later steps may appear incorrect only through dependence on an earlier corrupted state, and causal-inference variants \citep{MaCausalAttr250908682} assign step-level blame via performance-inversion and Shapley credit --- yet every member of the wave reads the trajectory's tool observations as faithful inputs to be reasoned over; propagated \emph{agent} error over a faithful log is exactly the case they handle, and a lying log is exactly the case they do not. All of these works treat the trajectory log (or, for Causal Agent Replay, the live re-executed trajectory) as faithful evidence once corrected for confounds of their own choosing --- Who\&When's decisive-error definition and MAST's taxonomy are both exercised on naturally occurring agent mistakes, AgenTracer's fault universe is agent actions, not tool returns, and Causal Agent Replay's validation universe is synthetic and mocked --- so none controls for, or even represents, the specific case this paper measures: a real attribution method's behavior on a real benchmark trajectory whose tool observation is corrupted yet well-formed.

\textbf{Security and reliability of the tool-return channel.} ChatInject \citep{ChatInjectWVhgFSKniL} constructs malicious tool outputs that forge chat-template role tags so that low-privilege tool content is interpreted as high-privilege instructions. MCP-SafetyBench \citep{MCPSafetyBench7XYjeL46co} embeds a multi-layer attack taxonomy, including tool poisoning, into executable tasks against real MCP servers. ResiliBench \citep{ResiliBenchKTZ56LG7jZ} builds a simulated tool registry with parameterized probabilistic error modes --- timeout, invalid input, operational failure --- to measure whether agents adapt workflows under execution uncertainty. All three establish that the tool-return channel is either actively attackable or pervasively unreliable, but none asks whether a downstream diagnosis correctly localizes the resulting failure: ChatInject and MCP-SafetyBench measure attack success on instruction-bearing payloads, a different threat model from a plainly wrong value in an ordinary field, and ResiliBench's error modes are overt by design, surfacing as a visible failure the agent must adapt to rather than a silent, well-formed lie. Nearest on the manipulated variable itself, Don't Blindly Trust It \citep{DontBlindlyTrust260621409} contrasts faithful, misleading, and absent tool feedback in a matched agent loop and finds a value inversion --- agents that benefit from clean tools can perform worse under misleading feedback than with no feedback at all --- establishing the phenomenon's cost to \emph{task performance} while leaving its cost to \emph{diagnosis} unmeasured, which is precisely the estimand here.

\textbf{Corrupted evidence beyond the tool return.} The memory-poisoning literature supplies the nearest neighbor to our estimand family. The Misattribution Gap \citep{MisattributionGap260522842} shows that memory-layer attacks --- a policy-formatted document laundered into a shared vector store and later resurfacing as trusted context --- produce agent misconduct that attribution and safety tooling systematically misread as model failure: across 64 documented failures, attribution consistently blamed the model, and four safety classifiers produced zero detections; the authors' counterfactual composition test then identifies the causal entry, the same conceptual move as our twin-replay oracle. Misattribution under a corrupted evidence channel is therefore no longer an untested idea, and we concede that priority plainly. The delta that defines this paper is threefold: the corrupted channel (a transient, per-step tool observation rather than a persistent memory store), the threat model (a non-adversarial, schema-valid wrong value in an ordinary field rather than an attack requiring provenance loss), and the diagnosis under test (the deployed plan/tool/execution attribution families --- LLM judges, trained tracers, runtime monitors --- run blinded with matched channel controls and a direction estimand, rather than safety classifiers and defender remediation choices).

\textbf{Runtime oversight and confound isolation.} SupervisorAgent \citep{SupervisorAgentpzFhtpkabh} adds a lightweight meta-agent that supervises high-risk interaction channels at runtime, activated by an LLM-free trigger filter keyed to explicit error occurrence, repetitive behavior, or excessive observation length. BreakingAgentBackbones \citep{BreakingAgentBackboneskga18ld70t} argues that end-to-end agent benchmarks make it difficult to attribute failures to the backbone LLM rather than to tool wrappers or control flow, and resolves this by isolating single LLM calls with their full context for security scoring rather than testing whether trajectory-level attribution mislocalizes a fault. The former's oversight signals are exactly what a schema-valid corruption does not emit; the latter motivates the attribution-confounding problem this paper measures but answers it by removing the trajectory rather than testing attribution over it.

\section{Conclusion}

This work reframes an assumption implicit in every automated failure-attribution pipeline --- that the trajectory log is trustworthy evidence --- as an empirical question with a precise, checkable structure, rather than treating a well-formed tool observation as self-evidently reliable. The central methodological wager is that a claimed attribution bias is only credible once causal decisiveness and surface undetectability are verified as separate properties of the injected fault, and once any bias is shown to be specific to the tool-observation channel rather than to the general difficulty of noticing a subtly wrong value; weaker designs risk mistaking incidental correlation, mapping noise, or generic value-spotting difficulty for a genuine attribution defect. Whichever pattern current methods exhibit, the finding is decision-relevant: a positive, direction-specific $\Delta$-trust gap would mean teams retraining planners or hardening execution in response to attribution verdicts are often chasing the wrong component, while a null result would mean the trace-as-evidence assumption is more robust than the structural argument suggests. Either outcome also completes a picture the memory-poisoning literature has begun to sketch for persistent stores \citep{MisattributionGap260522842} --- poisoned memory misread as model failure --- by measuring the corrupted-evidence channel that every attribution method reads on every step, under the threat model that requires no attacker at all. Two limitations bound this claim from the outset: the study covers single-observation corruptions verified one at a time, leaving compounding multi-observation corruption for follow-up work; and the mapping rubric's reliability (RQ7) may itself surface disagreement on trajectories where the plan/tool/execution boundary is genuinely ambiguous --- a limitation this design reports rather than resolves.

\vspace{1em}
\hrule
\vspace{1em}

\section*{Pre-handoff checklist}
\begin{itemize}[leftmargin=1.4em]
\item \textbf{Organization/navigation} --- narrative flows problem $\to$ method $\to$ experiments $\to$ related work without forward references to undefined terms.
\item \textbf{Terminology discipline} --- ``schema-valid corruption,'' ``decisive cause,'' ``misattribution rate,'' ``$\Delta$-trust gap,'' ``mapping rubric'' each defined at first use in Problem Definition or Methodology, used consistently thereafter; no internal stage IDs (\texttt{sp\_*}, \texttt{m\_*}, \texttt{rq\_*}) leaked into prose --- checked directly against the discussion pipeline log.
\item \textbf{Notation with worked example} --- the formal setup section defines $(\tau, q, a_t, o_t, V, \mu, \mathrm{replay}, A, \hat{\ell}, R_A)$ and the Problem Definition/Methodology sections consistently reuse this notation; no separate worked numeric example is included since no numbers exist yet (would require fabricating results).
\item \textbf{RQ framing matches outcomes} --- all 7 RQs are graded/quantitative (rates, proportions, $\kappa$), each answered in its own subsection with an explicit falsifiable expected outcome, not a binary yes/no.
\item \textbf{No typesetting artifacts} --- only explicit placeholder markers in Results, per the drafting convention this is a pre-registered spec, not a report.
\item \textbf{Artifact plan} --- the corruption injector, twin-replay oracle, blinded harness, and mapping rubrics are the natural release artifacts; not yet built, so not claimed as available.
\end{itemize}

\section*{Hand-off}
This is a draft, not a finished paper --- experiments have not run. Natural next steps: prose polish and venue-specific LaTeX formatting once a target venue is chosen, or revising an earlier stage if piloting surfaces a problem (e.g., if verified yield in RQ2 is much lower than expected).

\bibliographystyle{plainnat}
\bibliography{refs}

\end{document}
